From Corollary \ref{det-poly-steps-in-ppoly} and Theorem \ref{det-lower-bound}, we know that $\CoT[\poly(n), \poly(n)] = \ppoly$.
In the same way, from Corollary~\ref{prob-poly-steps-in-rppoly} and Theorem \ref{prob-lower-bound}, it follows that $\RCoT[\poly(n), \poly(n)] = \Rppoly$.
Thus, we obtain the following theorem.
%
\begin{theorem}\label{ppoly-rppoly-iff-cot-rcot}
    $\ppoly = \Rppoly \iff \CoT[\poly(n), \poly(n)] = \RCoT[\poly(n), \poly(n)]$.
\end{theorem}
%
This means that randomness boosts the power of transformers if and only if it also does for polynomial-size circuits, which is not known, as probabilistic circuits have not been deeply studied.

\bigskip

A more interesting question arises when we require the transformers to be efficiently constructible.

Observe that the proof of Theorem \ref{det-lower-bound} yields a polynomial algorithm that, for each circuit, constructs its corresponding deterministic transformer.
This motivates the definition of $\UCoT[T(n), d(n)]$, which is $\CoTtndn$ restricted to families of transformers such that there exists a deterministic Turing machine that runs in polynomial time and, given the input length $1^{n}$, outputs the transformer $\TF_n$.
Since $\ppoly$ restricted to uniform circuits coincides with $\Ptime$ \rafa{agregar ref a \ref{uppoly-equals-p}}, we obtain the following corollary.
%
\begin{corollary}\label{p-UCoT}
    $\Ptime \subseteq \UCoT[\poly(n), \poly(n)]$
\end{corollary}
%
Observe that Corollary \ref{p-UCoT} and Corollary \ref{det-poly-steps-in-ppoly} together imply Theorem \ref{p-equals-ucot}.
%
\begin{theorem}\label{p-equals-ucot}
    $\Ptime = \UCoT[\poly(n), \poly(n)]$
\end{theorem}


\medskip


These results have their analogous in the probabilistic setting.
In the same way as in the proof of Theorem \ref{det-lower-bound}, the proof of Theorem \ref{prob-lower-bound} yields a polynomial algorithm that, for each randomized circuit, constructs its corresponding probabilistic transformer.
Next, we define $\URCoT[T(n), d(n)]$, which is $\RCoTtndn$ restricted to families of transformers such that there exists a deterministic Turing machine that runs in polynomial time and, given the input length $1^{n}$, outputs the transformer $\TF_n$.
Since $\Rppoly$ restricted to uniform circuits coincides with $\BPP$ \rafa{agregar ref a \ref{urppoly-equals-bpp}}, we obtain the following corollary.
%
\begin{corollary}\label{p-URCoT}
    $\BPP \subseteq \URCoT[\poly(n), \poly(n)]$
\end{corollary}
%
Observe that Corollaries~\ref{p-URCoT} and~\ref{prob-poly-steps-in-rppoly} together imply the following theorem.
%
\begin{theorem}\label{bpp-equals-urcot}
    $\BPP = \URCoT[\poly(n), \poly(n)]$
\end{theorem}

\bigskip

Therefore, we deduce the uniform version of Theorem \ref{ppoly-rppoly-iff-cot-rcot}.
%
\begin{theorem}\label{p-bpp-iff-ucot-urcot}
    $\Ptime = \BPP \iff \CoT[\poly(n), \poly(n)] = \RCoT[\poly(n), \poly(n)]$.
\end{theorem}

Although $\Ptime = \BPP$ is an open question in complexity theory, this theorem is an interesting result.
It implies that, when restricted to uniform families of transformers, the constrained access to randomness that transformers have is equivalent in expressive power to the unrestricted access allowed to probabilistic polynomial-time Turing machines.